\documentclass[a4paper, 10pt]{scrartcl}


\usepackage{fontspec}
\newcommand{\mainFont}{NotoSerif}
\setmainfont{\mainFont}[
  Path           = ./fonts/\mainFont/static/,
  Extension      = .ttf,
  UprightFont    = *-Regular,
  BoldFont       = *-Bold,
  ItalicFont     = *-Italic,
  BoldItalicFont = *-BoldItalic,
]

\usepackage{xcolor}
\usepackage{tikz}
%\usepackage{lmodern} % Allows arbitrary font scaling without size limits
\usepackage{graphicx}
\usepackage{fontawesome5} % Icons: \faEnvelope, \faPhone, \faGithub, \faLinkedin, \faMapMarker, ...
\usepackage{hyperref} % Clickable mailto:/tel:/http links

\usetikzlibrary{calc} % Required for efficient coordinate calculations

\definecolor{background1}{HTML}{00171f}
\definecolor{background2}{HTML}{00335a}
\definecolor{col3}{HTML}{007da9}
\definecolor{picBorderCol}{HTML}{ffffff}
\definecolor{linkCol}{HTML}{490a83}

% Make links colored instead of the default red/boxed look
\hypersetup{
	colorlinks=true,
	urlcolor=linkCol,
	linkcolor=linkCol,
}

% options for text: \small, \normalsize, \large, \Large, \LARGE, \huge, \Huge

\pagestyle{empty}
\begin{document}

\hyphenpenalty=10000
\exhyphenpenalty=10000

\begin{tikzpicture}[remember picture, overlay, shift={(current page.north west)}]
	% (0,0) is now the top-left corner.




	% --- background ------------------------------------
	\pgfmathsetmacro{\splitpointX}{4}
	\pgfmathsetmacro{\splitpointY}{-2.8}
	\pgfmathsetmacro{\cornerpointUpX}{7.5}
	\pgfmathsetmacro{\cornerpointUpY}{-4.5}

	%\pgfmathsetmacro{\cornerpointDownX}{\cornerpointUpX}
	%\pgfmathsetmacro{\cornerpointDownY}{-7.8}

	\pgfmathsetmacro{\lastpointX}{4}
	\pgfmathsetmacro{\lastpointY}{-6.2}

	\fill[col3]
        (-1,1) --
        (-1,\lastpointY) --
        % Bézier S-curve connecting the bottom level to the top level
        (\lastpointX,\lastpointY) .. controls (\lastpointX + 2, \lastpointY) and (\cornerpointUpX - 2, \cornerpointUpY) .. (\cornerpointUpX,\cornerpointUpY) --
        (22,\cornerpointUpY) --
        (22,1) --
        cycle;


	%\fill[background2]
        %(-1,1) --
        %(-1,\splitpointY) --
        %% Bézier S-curve connecting the bottom level to the top level
        %(\splitpointX,\splitpointY) .. controls (\splitpointX + 2, \splitpointY) and (\cornerpointDownX - 2, \cornerpointDownY) .. (\cornerpointDownX,\cornerpointDownY) --
        %(22,\cornerpointDownY) --
        %(22,1) --
        %cycle;

	\fill[background1]
        (-1,1) --
        (-1,\splitpointY) --
        % Bézier S-curve connecting the bottom level to the top level
        (\splitpointX,\splitpointY) .. controls (\splitpointX + 2, \splitpointY) and (\cornerpointUpX - 2, \cornerpointUpY) .. (\cornerpointUpX,\cornerpointUpY) --
        (22,\cornerpointUpY) --
        (22,1) --
        cycle;













	% --- title name ------------------------------------
	\node[anchor=north west, text=white, font=\fontsize{42}{50}\selectfont\bfseries] at (7, -2.0)
	{
		Lucas Furer
	};







	% --- Contact section ------------------------------------
	% \faXXX gives the icon, \href{...}{...} makes it clickable.
	% mailto: opens the default mail client with the address
	% pre-filled; tel: does the same for a phone/dialer app.
	% On both PC and phone PDF readers you can usually also
	% right-click (or long-press) the link to copy it directly.
	\node[anchor=north west, align=left] at (0.5, -2.92)
	{
		\renewcommand{\arraystretch}{1.4}
		\begingroup
		\hypersetup{urlcolor=black}
		\begin{tabular}{@{}l l@{}}
			\faLinkedin   & \href{https://www.linkedin.com/in/lucas-furer-22a38a323/}{lucas-furer} \\
			\faGithub     & \href{https://github.com/LucasFurer}{LucasFurer} \\
			\faEnvelope   & \href{mailto:lucas.furer@gmail.com}{lucas.furer@gmail.com} \\
			\faPhone      & \href{tel:+31652772875}{+31 6 5277 2875} \\
			\faMapMarker  & Delft, NL \\
		\end{tabular}
		\endgroup
	};








	% --- section variables ------------------------------------
	\coordinate (sectStart) at ($(3.0,-7.2)$);


	\pgfmathsetmacro{\horGapSmall}{0.15} % space between text and the line
	\pgfmathsetmacro{\horGapLarge}{16.0} % space between line and dates
	\pgfmathsetmacro{\textWidth}{15.0} % width of the text

	\pgfmathsetmacro{\vertTextGap}{0.05} % spacing between text
	\pgfmathsetmacro{\vertSubGap}{0.1} % spacing between sub headers
	\pgfmathsetmacro{\vertSectGap}{0.3} % spacing between sections






	% --- text ------------------------------------
	\coordinate (eduStart) at ($(sectStart)$);


	%\node[anchor=north east, align=left, font=\Large] 
	%	at ($(eduStart) + (-\horGapSmall,0)$)
	%{
	%	\textbf{Education}
	%};


	\node[anchor=north west, align=left, text width=\textWidth cm] (eduEntryA) 
		at ($(eduStart) + (\horGapSmall,0)$)
	{
Dear Optiver Recruitment Team,\\
\vspace{1em}
I am writing to apply for the Graduate Software Engineer position at Optiver in Amsterdam. I recently graduated from TU Delft with an MSc in Computer Science, and I am eager to apply my background in C++, high-performance computing, and algorithmic problem-solving to the technically demanding environment of financial markets.\\
\vspace{1em}
What attracts me to Optiver is the opportunity to work on mission-critical systems where speed, reliability, and precision have a direct impact. I am interested in the financial industry and excited by the prospect of working closely with traders, researchers, and engineers to solve complex problems in a fast-moving environment. The combination of challenging technical work, continuous improvement, and close collaboration strongly appeals to me.\\
\vspace{1em}
During my Master's thesis, I implemented a fast multipole N-body solver and integrated it into t-SNE, with the core work developed in C++. This project required me to work with computationally intensive algorithms and consider performance when designing and implementing the solution. My Bachelor's thesis also involved developing a C++ implementation to study the behaviour of ray reflection distributions under microstructure smoothing. Through these projects, I developed a strong interest in tackling technically complex problems and building efficient, well-structured software.\\
\vspace{1em}
Alongside my academic work, I have gained experience working collaboratively in software development. During my internship at Feedback Analytics, I led a five-person team and worked directly with the client to gather requirements and deliver a visualization tool. This experience strengthened my ability to communicate clearly, take ownership, and collaborate effectively while keeping a project moving forward. I have also worked with Java, Python, SQL, Linux, Git, and CI/CD tooling, giving me experience across different areas of the software development process.\\
\vspace{1em}
I am an eager learner who enjoys exploring new technologies and challenging assumptions to find better solutions. Optiver's emphasis on continuous learning, engineering excellence, and using cutting-edge technology to solve problems in data-rich financial markets makes this opportunity particularly exciting to me. I would welcome the opportunity to learn from and contribute to Optiver's engineers, traders, and researchers while continuing to develop as a software engineer.\\
\vspace{1em}
Thank you for considering my application. I would welcome the opportunity to discuss how my technical background, analytical mindset, and interest in the financial industry could contribute to Optiver.\\
\vspace{1em}
Sincerely,\\
Lucas Furer
	};
	%\node[anchor=north east, align=left] 
	%	at ($(eduEntryA.north west) + (\horGapLarge,0)$)
	%{
	%	\textbf{2023-2026}
	%};



	\coordinate (eduEnd) at ($(eduStart |- eduEntryA.south)$);
	\draw[background2, line width=1pt]
		($(eduStart)$) --
		($(eduEnd)$);

	




















\end{tikzpicture}

\end{document}
